\section{Complexidade e porte: os grandes divisores da rede}
\label{sec:complexidadeporte}

Antes de qualquer comparação entre modelos de gestão, é preciso
reconhecer o quanto a complexidade e o porte, sozinhos, separam os
hospitais. As Figuras~\ref{fig:faixaA} a~\ref{fig:faixaC} mostram a
mediana de cada indicador ao longo das faixas Barcelona de complexidade
(da faixa 2, menor, à 6, maior), e as Figuras~\ref{fig:porteA}
a~\ref{fig:porteC}, ao longo dos portes.

O padrão é de escada: o custo mediano por saída sobe de \textbf{R\$~495
na faixa 2 para R\$~2.772 na faixa 6} (mais de cinco vezes), o tempo de
permanência vai de 3,27 para 6,30 dias, a fração de alta complexidade
sai de zero e chega a 22,5\%, e a ocupação de internação sobe de 37,5\%
para a casa dos 90\%. A mortalidade cresce com a complexidade até a
faixa 5 (de 2,88\% para 6,40\%) e \textbf{recua na faixa 6 (5,43\%)},
um sinal de que os hospitais de ponta convertem estrutura em resultado.
Vale notar que metade dos 49 hospitais da faixa 2 não opera UTI. No
porte, a mesma escada: custo mediano de \textbf{R\$~834 nos médios,
R\$~1.470 nos grandes e R\$~2.754 nos especiais}.

A conclusão para todo o restante do documento: \textbf{complexidade e
porte precisam ser sempre levados em conta ao comparar categorias de
gestão}, sob pena de atribuir ao modelo de gestão diferenças que são,
na verdade, de perfil assistencial. E, pela salvaguarda definida na
Seção~\ref{sec:painel}, a complexidade entra nas análises de
mortalidade apenas pela versão estrutural do escore.

\begin{figure}[H]
\centering
\subfig{fig_ae_05_faixa_mort_all.png}{Mortalidade geral}
\subfig{fig_ae_05_faixa_tmp.png}{TMP}
\caption{Mediana da mortalidade e do TMP por faixa Barcelona de
complexidade (2 a 6), com o valor sobre cada barra.}
\label{fig:faixaA}
\end{figure}

\comoler{cada barra é a mediana do indicador naquela faixa de
complexidade, com o valor impresso em cima. Na mortalidade, repare no
degrau que sobe até a faixa 5 e desce na 6; no TMP, a escada é sempre
ascendente.}

\begin{figure}[H]
\centering
\subfig{fig_ae_05_faixa_custo_saida.png}{Custo por saída}
\subfig{fig_ae_05_faixa_pct_alta_complex.png}{Fração alta complexidade}
\caption{Mediana do custo por saída e da fração de alta complexidade
por faixa Barcelona.}
\label{fig:faixaB}
\end{figure}

\comoler{a escada do custo é a mais íngreme do documento: da faixa 2
para a 6, a mediana mais que quintuplica. A alta complexidade parte do
zero absoluto nas faixas baixas.}

\begin{figure}[H]
\centering
\subfig{fig_ae_05_faixa_ocupacao_internacao.png}{Ocupação de
internação}
\subfig{fig_ae_05_faixa_ocupacao_uti.png}{Ocupação de UTI}
\caption{Mediana das taxas de ocupação por faixa Barcelona.}
\label{fig:faixaC}
\end{figure}

\comoler{quanto mais complexo o hospital, mais cheio ele opera. Na
UTI, a barra da faixa 2 vale zero porque metade daqueles hospitais não
tem UTI.}

\begin{figure}[H]
\centering
\subfig{fig_ae_06_porte_mort_all.png}{Mortalidade geral}
\subfig{fig_ae_06_porte_tmp.png}{TMP}
\caption{Mediana da mortalidade e do TMP por porte (classificação fixa
pela mediana de leitos: 173 hospitais médios, 130 grandes e 11
especiais).}
\label{fig:porteA}
\end{figure}

\comoler{mesma leitura das figuras por faixa, agora por porte:
hospitais maiores seguram o paciente por mais tempo e têm mortalidade
um pouco maior, coerente com casos mais graves.}

\begin{figure}[H]
\centering
\subfig{fig_ae_06_porte_custo_saida.png}{Custo por saída}
\subfig{fig_ae_06_porte_pct_alta_complex.png}{Fração alta complexidade}
\caption{Mediana do custo por saída e da fração de alta complexidade
por porte.}
\label{fig:porteB}
\end{figure}

\comoler{o custo mediano dos hospitais especiais é mais de três vezes o
dos médios, e a alta complexidade se concentra quase toda nos portes
grande e especial.}

\begin{figure}[H]
\centering
\subfig{fig_ae_06_porte_ocupacao_internacao.png}{Ocupação de
internação}
\subfig{fig_ae_06_porte_ocupacao_uti.png}{Ocupação de UTI}
\caption{Mediana das taxas de ocupação por porte.}
\label{fig:porteC}
\end{figure}

\comoler{hospitais maiores operam mais cheios, tanto na internação
quanto na UTI: mais um retrato de que porte e pressão assistencial
andam juntos.}
